\section{Small stabilizers}
\label{small:section}

We use Proposition~\ref{cone:reduction} and
Lemmas~\ref{cone:approximation} and~\ref{cone:orbifold}.  Thus
\(C=\Spec S\) is a normal Gorenstein Fano cone of dimension \(n=N+1\),
where \(N\geq4\), its puncture \(C^\circ=C\setminus\{o\}\) is smooth,
and the polarized cone \((C,\xi)\) is K-polystable, with minimizing
Reeb vector normalized by \(A(\xi)=n\).  Fix homogeneous torus eigenfunctions
\(f_1,\ldots,f_q\) generating \(S\), with torus characters
\(\alpha_1,\ldots,\alpha_q\).  For a primitive positive integral grading
\(\zeta\), write
\[
 r=A(\zeta),\qquad
 B=[C^\circ/\Gm],\qquad -K_B=rL.
\]
Here \(L\) is the faithful ample orbifold line bundle of weight one.
The maximum stabilizer order for this grading is denoted by
\(m_{\max}\).  All constants in this section may depend on the fixed cone,
its torus and the displayed generators.

\begin{theorem}[Small-stabilizer classification]
\label{small:classification}
Let \(\zeta_i\) be primitive positive integral gradings satisfying
\begin{equation}
 \zeta_i=\frac{r_i}{n}\xi+\epsilon_i,
 \qquad r_i=A(\zeta_i),\qquad \|\epsilon_i\|\longrightarrow0,
 \qquad m_{\max}(\zeta_i)\leq\frac{r_i}{N}.
 \label{small:fast-sequence}
\end{equation}
For all sufficiently large \(i\), every stabilizer is trivial and
\(r_i\in\{N,N+1\}\).  More precisely, the polarized quotient is
\[
 (B_i,L_i)\cong
 \begin{cases}
  (\PP^N,\cO_{\PP^N}(1)),&r_i=N+1,\\
  (Q^N,\cO_{Q^N}(1)),&r_i=N,
 \end{cases}
\]
where \(Q^N\subset\PP^{N+1}\) is a smooth quadric.  Consequently the cone
\(C\) is respectively smooth or an ordinary double point.  The assertion
includes the constant exact primitive grading when the ray of \(\xi\)
is rational.
\end{theorem}

We begin with a minimal orbifold rational curve.  Its tangent characters
either force \(r=N\) and trivial stabilizers, or its deformations yield a
finite evaluation map and hence a nonzero symmetric tensor of negative
weight.  A metric lower bound for homogeneous symmetric tensors, together
with an approximation estimate uniform in their rank, gives the remaining
rigidity.

\subsection{A minimal orbifold rational curve}

For the moment fix one primitive grading satisfying
\begin{equation}
 m_{\max}\leq r/N.
 \label{small:stabilizer-assumption}
\end{equation}
The orbifold \(B\) is smooth and effective, is proper with projective
coarse space, and has ample anticanonical bundle \(rL\).  We use the
orbifold rational-curve theorem of Li--Zhou
\cite[Proposition~2.6]{LZ25}.  In this setting it supplies a
representable nonconstant map from \(\PP(1,\ell)\), where the only
possible stack point is infinity, whose anticanonical degree plus the
inverse age at that point is at most \(\dim B+1\).  In particular its
anticanonical degree is at most \(n\).  No positivity assumption on
\(T_B\) is required for this assertion.

Choose a map of least positive \(L\)-degree among all representable
nonconstant maps of this form, allowing \(\ell=1\).  Such a minimum
exists: representability embeds the source stabilizer in a target
stabilizer, so the denominators of these degrees divide one of the
finitely many integers at most \(m_{\max}\).  Write
\[
 f:\PP(1,\ell)\longrightarrow B,
 \qquad f^*L\cong\cO(e),\qquad e>0,
\]
and let \(m\) be the stabilizer order of its value at infinity.
This endpoint order need not equal \(m_{\max}\).  Thus
\(\ell\mid m\), and minimality and the existence theorem give
\(re/\ell\leq n\).  On the other hand,
\(r/m\geq N\) by Equation~\eqref{small:stabilizer-assumption}.  Hence
\begin{equation}
 1\leq\frac{em}{\ell}
 =\frac{re/\ell}{r/m}
 \leq\frac{N+1}{N}<2.
 \label{small:primitive-curve-degree}
\end{equation}
The expression \(em/\ell\) is an integer.  We obtain
\begin{equation}
 e=1,\qquad \ell=m,
 \qquad N\leq r/m\leq n.
 \label{small:curve-range}
\end{equation}

Choose an identification \(f^*L\cong\cO(1)\).  Passing to its grading
torsor gives an equivariant map
\[
 F:\A^2\setminus\{0\}\longrightarrow C^\circ,
 \qquad F(t x,t^m y)=t\cdot F(x,y).
\]
Since \(\A^2\) is normal and \(C\) is affine, the coordinate functions
extend uniquely across the origin.  Thus \(F:\A^2(1,m)\to C\) is a
homogeneous polynomial map.  Set \(P=F(0,1)\); its stabilizer is
\(\mu_m\).  Conversely, a homogeneous map with no preimage of the
vertex outside the origin, and with this degree-one torsor
identification, gives a representable map to \(B\): its maps on source
inertia groups are the inclusions inherited from the identity
homomorphism of the grading groups.

\begin{lemma}[Deformations with a fixed lift]
\label{small:deformation}
Let \(b_1',\ldots,b_n'\in\{1,\ldots,m\}\) represent the characters of
\(\mu_m\) on \(T_{C,P}\), with the trivial character represented by
\(m\).  Every component through \(F\) of the scheme of homogeneous maps
\(F':\A^2(1,m)\to C\) satisfying \(F'(0,1)=P\) has dimension at least
\begin{equation}
 \chi=n+\frac rm-\sum_{j=1}^n\frac{b_j'}m.
 \label{small:chi}
\end{equation}
In particular this dimension is at least \(n\), except possibly when
\(r/m=N\) and the entire tangent action at \(P\) is trivial.
\end{lemma}

\begin{proof}
Homogeneous maps form an affine scheme of finite type: in fixed
homogeneous generators of \(S\), each coordinate of the map is a
polynomial of a specified weighted degree, and the relations of \(S\)
give polynomial equations in their finitely many coefficients.  The
condition of avoiding the vertex outside the origin is open, as may be
checked on the proper weighted line.  Near the given map, deformation
theory can therefore be performed entirely in the smooth target
\(C^\circ\).

Let \(E\) be the rank-\(n\) bundle on \(\PP(1,m)\) obtained by
descending \(F^*T_{C^\circ}\).  Its determinant is \(\cO(r)\), so
\(\deg E=r/m\).  We are deforming an equivariant map with its source
and its grading torsor fixed; equivalently, we are deforming a section
of the associated smooth bundle, with the canonical section obstruction
theory of \cite[Theorem~4.13 and Equation~(4.18)]{Webb22}.
There is no quotient by source
automorphisms or by changes of the torsor identification.  If
\(D_\infty\) denotes the residual gerbe at infinity, the tangent and
obstruction spaces for fixing its value are
\[
 H^0(\PP(1,m),E(-D_\infty)),\qquad
 H^1(\PP(1,m),E(-D_\infty)).
\]
Thus their Euler characteristic is a lower bound for the dimension of
every component through the map.

For completeness, orbifold Riemann--Roch here says
\(\chi(E)=n+\deg E-\operatorname{age}(E_\infty)\).
It follows by taking the coarse pushforward: locally diagonalizing at
the stack point subtracts the fractional character weights from the
degree, after which ordinary Riemann--Roch on \(\PP^1\) applies.
Restriction to \(D_\infty\) has Euler characteristic equal to the
dimension of the invariant fiber, because \(\mu_m\) is linearly
reductive.  Subtracting that invariant dimension is precisely the
effect of representing each trivial character by \(m\) instead of
zero.  The exact sequence for restriction to \(D_\infty\) therefore
gives Equation~\eqref{small:chi}.

Since \(\sum b_j'/m\leq n\), we have \(\chi\geq r/m\).  The Euler
characteristic \(\chi\) is an integer.  If \(r/m>N\), it is at least
\(N+1=n\).  If \(r/m=N\)
and some character is nontrivial, the inequality is strict and the
same conclusion follows.
\end{proof}

\subsection{Finite evaluation and its symmetric tensor}

Assume that the exception in Lemma~\ref{small:deformation} does not
occur.  Let \(H\) be the reduced closure of an irreducible component
of the open locus of maps avoiding the vertex outside the origin,
in the affine scheme of homogeneous maps with \(F'(0,1)=P\).
Then \(H\) is integral and
\begin{equation}
 \dim H\geq n.
 \label{small:H-dimension}
\end{equation}
Precomposition with \((x,y)\mapsto(a x,y)\) preserves the chosen
component because the acting group \(\Gm\) is connected.  It gives a
positive grading on \(H\): every coefficient involving \(x\) has positive weight,
whereas the coefficients involving only \(y\) are fixed by \(P\).
All maps contract to the same map obtained by deleting the terms
involving \(x\).  Consider the equivariant evaluation morphism
\begin{equation}
 \operatorname{ev}:H\longrightarrow C,
 \qquad F'\longmapsto F'(1,0).
 \label{small:evaluation}
\end{equation}

\begin{lemma}[Finite evaluation]
\label{small:finite-evaluation}
The morphism in Equation~\eqref{small:evaluation} is finite and
surjective.  Its fiber over \(o\) is supported at the unique contracted
map.
\end{lemma}

\begin{proof}
Take \(F_0\in H\) with \(F_0(1,0)=o\).  Choose a pointed curve in
\(H\) through \(F_0\), with general point in the locus avoiding the
vertex, and normalize it.  Its generic map gives a twisted stable
graph map to \(B\times\PP^1\): the second map is the coarse
identification of \(\PP(1,m)\) with \(\PP^1\), and infinity is marked.
This graph is stable even before any contraction because its second
map has degree one.  After a finite base change, properness of twisted
stable maps gives a limiting graph map
\cite[Theorem~4.3]{AOV11}.  The target is a tame proper smooth
Deligne--Mumford stack with projective coarse space, as required by
that theorem.

There is a unique component of the limiting coarse source of second
degree one.  It maps isomorphically to \(\PP^1\); call it the
horizontal component.  All other components form rational trees
attached to it, since the total arithmetic genus is zero.
The first map on the horizontal component is the coarse rational
map defined by \(F_0\), after cancellation of its basepoints.
Indeed the two maps agree on the open set where the limiting
homogeneous coordinates do not vanish simultaneously.

We spell out why a positive-degree tail occurs over the finite
basepoint \([1:0]\).  Choose a sufficiently divisible integer \(k\)
such that \(kL\) descends to a very ample bundle on the coarse space
of \(B\), and \(m\mid k\).  Its pullback along the original maps is
the coarse bundle \(\cO_{\PP^1}(k/m)\).  The limiting projective
coordinates are binary forms of that degree.  They have a common
zero at \([1:0]\), because \(F_0(1,0)=o\), and are not all zero,
because \(F_0(0,1)=P\ne o\).  Cancelling their common divisor loses
positive degree at \([1:0]\).  A general hyperplane avoiding the
value of the canceled map at that point has the lost positive number
of inverse-image points specializing there.  In the limiting stable
map these points cannot lie on the horizontal component.  They are
therefore accounted for by positive first degree on a tree over
\([1:0]\).  This also proves the degree assertion directly, without
making an assumption about components that might occur over
infinity.

Choose a terminal leaf of that tree.  It does not contain the
infinity marking, since its second image is \([1:0]\).  It has one
node and no markings.  If its first map had degree zero, it would be
constant in both factors and unstable: stability requires the induced
coarse map to be a Kontsevich stable map
\cite[Section~4]{AOV11}, and an unmarked constant terminal leaf has
only one special point. Hence its first degree is
positive.  Its normalization is a twisted rational curve with at
most one stack point: the only allowed stack points of a twisted
source are nodes and markings.  Projection from the product target
to \(B\) preserves representability, since \(\PP^1\) has trivial
inertia.  The leaf consequently gives a representable nonconstant
map \(\PP(1,\ell')\to B\) of the type used in the minimum.

Its \(L\)-degree is at least the original minimal degree \(1/m\).
Degree is additive on the limiting source and every component has
nonnegative \(L\)-degree, while the total degree is \(1/m\).
Thus this leaf consumes the full first degree.  The horizontal
component has first degree zero, and its constant value is the
coarse image of \(P\).  It follows that the image of the affine map
\(F_0\) lies in the closure of the grading orbit of \(P\).

The normalization of that orbit closure is \(\A^1\), with a
parameter of grading weight \(m\), normalized to have value one at
\(P\).  To see the weight, the gcd of the degrees of the coordinates
nonzero at \(P\) is its stabilizer order \(m\).  The map from the
normal source \(\A^2\) factors through this normalization: its
parameter is integral over the coordinate ring of the orbit curve
and hence regular on \(\A^2\).  The parameter pulled back along
\(F_0\) is a weighted homogeneous polynomial of degree \(m\), so it
has the form
\[
 y+c x^m.
\]
The coefficient of \(y\) is one because the value at \((0,1)\) is
fixed.  Evaluation at \((1,0)\) is the vertex only if \(c=0\).
Therefore the vertex fiber is supported at the single contracted map.

Finally, let \(R_H\) be the positively graded coordinate ring of
\(H\).  The preceding fiber statement says that
\(R_H/S_+R_H\) is finite dimensional.  Lift a homogeneous basis of
this quotient.  Induction on the positive degree shows that these
lifts generate \(R_H\) as an \(S\)-module; this is graded
Nakayama.  Thus evaluation is finite.  Its image is closed and has
dimension \(\dim H\geq n\), by Equation~\eqref{small:H-dimension}.
Since \(C\) is integral of dimension \(n\), the image is all of
\(C\), and \(\dim H=n\).
\end{proof}

\begin{lemma}[A tensor of negative integral weight]
\label{small:norm-tensor}
Let \(b\) be the degree of the finite evaluation in
Lemma~\ref{small:finite-evaluation}.  There is a nonzero algebraic
section
\[
 s\in H^0(C^\circ,\Sym^bT_{C^\circ})
\]
of integral-grading weight \(-bm\).
\end{lemma}

\begin{proof}
The source automorphisms
\((x,y)\mapsto(x,y+t x^m)\), for \(t\in\CC\), preserve the
homogeneous-map scheme, the fixed value at \((0,1)\), and its open
locus avoiding the vertex.  The connected additive group therefore
preserves the component \(H\).  Its vector field is nonzero: a
coordinate nonzero at \(P\) has a nonzero pure power of \(y\), whose
derivative under this action is nonzero in characteristic zero.
Differentiating evaluation gives a regular vector \(V\) along
evaluation over \(C^\circ\).  It is generically nonzero, since the
finite dominant evaluation is generically separable and hence has
injective differential on a nonempty smooth open set.  For a target
coordinate of degree \(d\), its component is
\(\partial_yF'(1,0)\), of coefficient-grading degree \(d-m\).
Thus \(V\), as a vector along evaluation, has derivation weight
\(-m\).

Over the generic point, take the symmetric product of the vectors
on all \(b\) sheets.  This gives a nonzero element of
\(\Sym^bT_{C^\circ}\) over \(\CC(C)\); the product, unlike the
ordinary trace, cannot cancel.  It is invariant under permutation
of the sheets and so descends to that field.  It is regular as
well.  On a smooth affine open set of \(C^\circ\) trivializing the
tangent bundle, the components of \(V\) lie in a finite algebra
over the base ring.  The coefficients of their symmetric product
are therefore integral over the base ring and lie in its fraction
field.  Normality of the base makes them regular.  This argument
does not require \(H\) to be normal; alternatively one can first
use its finite normalization.  The construction is equivariant
and multiplies weights by \(b\), giving weight \(-bm\).
\end{proof}

\subsection{A bound for homogeneous symmetric tensors}

The tensor just constructed has negative weight for the integral grading.
To exploit this, we first bound the Reeb weight of a homogeneous
symmetric tensor.  By the Fano cone
Yau--Tian--Donaldson theorem
\cite[Theorem~1.1 and Definitions~2.1--2.2]{CS19}, \(C^\circ\) carries a
Ricci-flat K\"ahler cone metric with Reeb vector \(\xi\).  The theorem
applies here because the puncture is smooth and the singularity is
Q-Gorenstein and klt.  The condition called K-stability in that reference
allows equality only for an isomorphic central fiber, and is the
K-polystability condition used here.  With \(A(\xi)=n\), the real
homothetic dilations have infinitesimal generator \(-J\xi\), and their
pullback multiplies the metric by the square of the dilation factor.

\begin{lemma}[Tensor weight bound]
\label{small:tensor-bound}
Let \(b\geq1\), and let \(s\) be a nonzero homogeneous algebraic section
of \(\Sym^bT_{C^\circ}\), of \(\xi\)-weight \(\tau\), with vector fields
weighted as derivations.  Then
\[
 \tau\geq-b.
\]
\end{lemma}

\begin{proof}
The Chern connection induced on \(\Sym^bT_{C^\circ}\) has zero mean
curvature: contraction of the curvature on \(T_{C^\circ}\) is Ricci
curvature, and contraction commutes with the induced tensor and symmetric
power operations.  Put \(u=|s|^2\).  The holomorphic Bochner formula and
Cauchy--Schwarz give, with a consistent positive complex-Laplacian
normalization,
\[
 \Delta u=|\nabla^{1,0}s|^2\geq0,
 \qquad |\partial u|^2\leq u\,\Delta u.
\]
Consequently \(\log(u+a)\) is subharmonic for each \(a>0\), and so is
\((u+a)^{p/2}\) for every \(p>0\), by the chain rule applied to the
increasing convex exponential function.  Letting \(a\) decrease to zero
shows that \(|s|^p\) is subharmonic in the distributional sense, including
across the zero set of \(s\).

Let \(\rho\) be the metric radial coordinate and \(M=\{\rho=1\}\) its
compact smooth link.  A contravariant tensor of rank \(b\) and derivation
weight \(\tau\) has norm of radial degree \(\tau+b\).  If
\(\tau+b<0\), choose \(p>0\) such that
\[
 d=p(\tau+b)\in(2-2n,0).
\]
Then \(|s|^p=\rho^d\phi\) for a nonnegative continuous function
\(\phi\) on \(M\) with positive integral.  The cone-Laplacian formula,
valid also distributionally, is
\[
 \Delta(\rho^d\phi)
 =\rho^{d-2}\bigl(d(d+2n-2)\phi+\Delta_M\phi\bigr).
\]
Integrating the angular term over \(M\) gives zero.  Testing
subharmonicity against a nonnegative radial function supported in an
annulus would therefore imply
\(d(d+2n-2)\int_M\phi\geq0\), which is impossible.  Thus
\(\tau+b\geq0\).
\end{proof}

\subsection{Uniform control of characters and the conclusion}

The tensor rank in Lemma~\ref{small:norm-tensor} can vary with the
grading.  The following estimate makes the approximation error
independent of that rank.

\begin{lemma}[Character estimate]
\label{small:character-estimate}
There is a constant \(c>0\) with the following property for all
sufficiently close gradings in Equation~\eqref{small:fast-sequence}.
If a nonzero rank-\(b\) tensor on \(C^\circ\) has integral-grading
weight \(-bm\), with \(m>0\), then it has a nonzero torus
eigencomponent of character \(\chi\) satisfying
\[
 \|\chi\|\leq cb.
\]
Consequently, if the tensor is as in
Lemma~\ref{small:norm-tensor}, then
\begin{equation}
 m\leq\frac rn+c\|\epsilon\|.
 \label{small:m-bound}
\end{equation}
\end{lemma}

\begin{proof}
For the fixed eigen-generators \(f_j\), positivity of \(\xi\) and
absolute convergence of the error give constants \(c_1,c_2>0\)
such that their integral degrees \(d_j=\alpha_j(\zeta)\) satisfy
\begin{equation}
 c_1r\leq d_j\leq c_2r
 \quad\text{for every }j.
 \label{small:generator-weights}
\end{equation}
These bounds include the exact rational grading.

The differentials \(df_j\) span the cotangent bundle of \(C^\circ\).
Pairing a symmetric tensor with all ordered \(b\)-tuples of these
differentials embeds its space of sections into a finite tuple of
regular functions on \(C^\circ\).  Normality extends these functions
over \(o\).  This also shows that the torus action on sections is
locally finite, so the tensor has a nonzero torus eigencomponent;
every such component retains its specified \(\zeta\)-weight.

Let \(\chi\) be the character of a nonzero component, and choose a
nonzero pairing with
\(df_{j_1}\otimes\cdots\otimes df_{j_b}\).  Its function has
character
\(\beta=\chi+\alpha_{j_1}+\cdots+\alpha_{j_b}\) and integral degree
\[
 0\leq d=-bm+d_{j_1}+\cdots+d_{j_b}\leq c_2rb.
\]
It has a polynomial lift in the generators consisting of monomials
of character \(\beta\), by torus equivariance of the polynomial
presentation of \(S\).  Each such monomial contains at most
\(d/(c_1r)\leq(c_2/c_1)b\) factors, by
Equation~\eqref{small:generator-weights}.  As the generator
characters are fixed, this proves \(\|\beta\|=O(b)\), and then
\(\|\chi\|=O(b)\), uniformly in the tensor, its rank and the
grading.

Put \(\tau=\chi(\xi)\).  The eigencomponent satisfies
Lemma~\ref{small:tensor-bound}, so \(\tau\geq-b\).  On the other
hand,
\[
 -bm=\chi(\zeta)=\frac rn\tau+\chi(\epsilon)
 \geq-\frac rn b-cb\|\epsilon\|.
\]
Dividing by \(b\) proves Equation~\eqref{small:m-bound}.
\end{proof}

\begin{proof}[Proof of Theorem~\ref{small:classification}]
For each grading choose the minimal curve and its stabilizer order
\(m\) as above.  If the exceptional case in
Lemma~\ref{small:deformation} occurs, then \(r/m=N\) and the
\(\mu_m\)-action on \(T_{C,P}\) is trivial.  Faithfulness of the
transverse tangent action implies \(m=1\), hence \(r=N\).
Equation~\eqref{small:stabilizer-assumption} then forces
\(m_{\max}=1\).

For the remaining indices, Lemmas~\ref{small:finite-evaluation}
and~\ref{small:norm-tensor} apply.  The character estimate and
Equation~\eqref{small:curve-range} give
\[
 0\leq nm-r\leq nc\|\epsilon\|.
\]
The left side is an integer.  Since the error tends to zero, we
have \(nm=r\) for all sufficiently large indices.  For a rational
ray with its constant exact grading the error is zero, so this
equality holds directly.

Now the finite evaluation has \(\dim H=n\), whereas its deformation
dimension satisfies
\[
 n=\dim H\geq n+\frac rm-\sum_{j=1}^n\frac{b_j'}m
 =2n-\sum_{j=1}^n\frac{b_j'}m\geq n.
\]
All equalities must hold.  In particular every \(b_j'=m\), so the
tangent action is trivial.  Its faithfulness again gives \(m=1\),
and now \(r=n=N+1\).  Finally
\[
 m_{\max}\leq\frac{N+1}{N}<2
\]
forces \(m_{\max}=1\).

The quotient is therefore an ordinary smooth projective Fano
\(N\)-fold and \(L\) an actual ample line bundle.  The
Kobayashi--Ochiai theorem in its ample-line-bundle formulation
\cite{KO73}, also supplied by \cite[Theorems~1.1 and~1.3]{ADK08}, applied to
\(-K_B=rL\) with \(r=N+1\) or \(N\), gives precisely the two
polarized varieties stated in the theorem.  The polarization is
the standard one, as also follows by substituting into their
Picard groups.  By the section-ring identity from the cone setup,
\[
 S=\bigoplus_{j\geq0}H^0(B,jL).
\]
This is the polynomial ring for \((\PP^N,\cO(1))\), and the
homogeneous coordinate ring of a nondegenerate quadratic
hypersurface for \((Q^N,\cO(1))\).  Its affine cone is respectively
smooth or the \(n\)-dimensional ordinary double point.
\end{proof}
